\section{Recommended Reading}

\begin{readingbox}
\textbf{Foundation}
\begin{itemize}[nosep]
  \item Eugene Izhikevich, \emph{Dynamical Systems in Neuroscience}, for
        neuronal dynamics, excitability classes, oscillation, and rhythm
        generation in a biologically concrete mathematical style.
  \item Michael Brin and Garrett Stuck, \emph{Introduction to Dynamical
        Systems}, for fixed points, stability, bifurcations, and invariant-set
        thinking at a more general mathematical level.
  \item Steven Strogatz, \emph{Nonlinear Dynamics and Chaos}, for the most
        readable bridge from phase portraits and limit cycles to threshold-like
        qualitative change.
  \item Arkady Pikovsky, Michael Rosenblum, and J\"urgen Kurths,
        \emph{Synchronization}, for coupling, entrainment, and phase-locking.
\end{itemize}
\textbf{Modern Research}
\begin{itemize}[nosep]
  \item Nature Reviews Neuroscience (2022) on attractor and integrator
        networks in the brain, as the chapter's main medium-lane review anchor.
  \item \emph{Metastability demystified: the foundational past, the pragmatic
        present and the promising future} (Nature Reviews Neuroscience, 2024,
        medium), for flexible stability and reconfiguration.
  \item \emph{Falling asleep follows a predictable bifurcation dynamic}
        (Nature Neuroscience, 2025, strong in its own biological domain), for
        a concrete transition model that legitimizes threshold-sensitive
        thinking.
  \item \emph{Self-orthogonalizing attractor neural networks emerging from the
        free energy principle} (arXiv, 2025, medium), for sequence formation
        and self-organizing attractor structure.
  \item \emph{Koopman-Nemytskii Operator: A Linear Representation of Nonlinear
        Controlled Systems} (arXiv, 2025, medium), for operator and feedback
        viewpoints as an advanced extension.
\end{itemize}
\textbf{Frontier}
\begin{itemize}[nosep]
  \item Advanced operator and consciousness-architecture proposals should be
        read as formal upgrade paths or hypothesis generators, not as finished
        behavioural proof.
\end{itemize}
\end{readingbox}
